\begin{corollary}
If $Q$ is a quasi-order and $Q$ is a better-quasi-order, then the full
tagged hierarchy $\mathsf{PowerQ}(Q)$ equipped with the four-clause relation
$\mathsf{PowerRel}(r)$ is a better-quasi-order:
\[
\mathsf{Bqo}(r)\longrightarrow
\mathsf{Bqo}(\mathsf{PowerRel}(r)).
\]
The node type has arbitrary nonempty small branching, as in
\noderef{PowerQ}; it is not restricted to countable nodes.
\end{corollary}
\begin{proof}
Assume $\mathsf{Bqo}(r)$ and suppose, for contradiction, that
$\mathsf{PowerRel}(r)$ is not a better-quasi-order.  By the definition of
better-quasi-order in \noderef{Bqo}, there is a locally constant bad
multi-sequence
$h:\mathsf{IncSeq}\to\mathsf{PowerQ}(Q)$.

Apply the reflection proposition \noderef{PowerQReflection} to this $h$.
It produces a locally constant $g:\mathsf{IncSeq}\to Q$ which is bad for
$r$.  This contradicts $\mathsf{Bqo}(r)$, again using the defining
quantifiers in \noderef{Bqo}.  Therefore no locally constant bad
multi-sequence exists for $\mathsf{PowerRel}(r)$, which is exactly the
claimed conclusion.  The quasi-order hypothesis is the displayed
$\mathsf{IsPreorder}$ assumption, and the reflection construction uses the
full arbitrary-small-indexed hierarchy supplied by \noderef{PowerQ}.
\end{proof}
